% ============================================================================
%  CAPÍTULO I: INTRODUCCIÓN
% ============================================================================

\section{Introducción}

En toda red corporativa existe un requisito silencioso pero innegociable:
traducir los nombres que los usuarios recuerdan (\texttt{web.sudoers.lan}) a
las direcciones IP que las m\'aquinas entienden (\texttt{192.168.0.2}). Sin
este puente, ni el navegador de un empleado podr\'ia abrir la p\'agina del
proyecto, ni el servidor de correo podr\'ia entregar un mensaje, ni una
impresora de red ser\'ia localizable. El servicio que materializa este puente
es el \textbf{DNS} (\textit{Domain Name System}) \citep{rfc1035}, y sobre
\'el se apoya la segunda pieza: el \textbf{servidor web} que publica los
contenidos que los clientes consumen a diario.

En la infraestructura de la organizaci\'on, el servidor de correo, el portal,
la impresi\'on y el panel de gesti\'on no se alcanzan tecleando IPs: se
alcanzan por \textbf{nombre} (\texttt{web.sudoers.lan},
\texttt{print.sudoers.lan}, \texttt{portainer.sudoers.lan}). Este cap\'itulo
introduce la problem\'atica, el prop\'osito de la actividad y el marco
te\'orico m\'inimo necesario para comprender el desarrollo del trabajo.

\subsection{Problem\'atica}

La red del laboratorio es un \'unico servidor Debian que concentra todos los
servicios de la organizaci\'on (identidad, archivos, correo, web) sobre una
sola red plana \texttt{192.168.0.0/24}. En ese escenario se plantean tres
carencias:

\begin{itemize}
  \item \textbf{Resoluci\'on de nombres ausente o manual:} sin un servidor
        DNS propio, cada cliente debe memorizar direcciones IP o depender de
        un archivo \texttt{/etc/hosts} local, esquema que no escala y que
        impide que un mismo nombre sea resuelto de forma consistente en toda
        la red.
  \item \textbf{Acceso a los servicios sin nombres amigables:} los recursos
        de red solo se alcanzan por IP, lo que obliga a los usuarios a
        recordar numeraciones y dificulta la administraci\'on y el crecimiento
        de la infraestructura.
  \item \textbf{Ausencia de una puerta \'unica y segura hacia la web:} sin un
        servidor web que concentre y publique los servicios por nombre, cada
        aplicaci\'on tendr\'ia que exponer su propio puerto, sin una capa de
        cifrado ni un punto \'unico de control.
\end{itemize}

La soluci\'on a estas carencias consiste en aprovechar un \'unico servidor
\textbf{Debian 12} que ya funciona como Controlador de Dominio y desplegar
sobre \'el dos piezas complementarias: el \textbf{DNS autoritativo} del
dominio, servido por el propio \textbf{Samba 4} (backend
\textbf{SAMBA\_INTERNAL}), y un \textbf{servidor web} basado en \textbf{Nginx}
que act\'ua como \textit{proxy inverso} de toda la intranet y que se publica
por nombre de dominio con \textbf{HTTPS}.

\subsection{Prop\'osito de la actividad}

La Actividad 8 tiene como prop\'osito configurar, sobre GNU/Linux con la
distribuci\'on \textbf{Debian 12}, el \textbf{DNS del dominio}
\texttt{sudoers.lan} administrado por el Controlador de Dominio de
\textbf{Samba} (zona autoritativa directa, registros de descubrimiento SRV y
reenv\'io hacia Internet) y el \textbf{servidor web} de la organizaci\'on
basado en \textbf{Nginx}, desplegado como contenedor Docker, que publique el
portal del proyecto y act\'ue como \textit{proxy inverso} de los servicios
publicados por nombre, con cifrado \textbf{HTTPS} mediante una CA interna.

Con ello se busca reproducir, en un entorno acad\'emico, la arquitectura de
resoluci\'on de nombres y de publicaci\'on de contenido que sostiene a una
organizaci\'on real.

\subsection{Marco te\'orico}

\subsubsection{El sistema de nombres de dominio (DNS)}

El DNS es un servicio distribuido y jer\'arquico que traduce nombres de
dominio a direcciones IP y viceversa \citep{rfc1034}. En un dominio de
Active Directory (AD), el DNS no es un accesorio: \textbf{es el transporte del
directorio}. Los clientes lo consultan para descubrir d\'onde est\'a el LDAP,
el Kerberos o el propio Controlador de Dominio mediante registros
\textbf{SRV} \citep{rfc2782}. Por eso un AD DC debe ser autoritativo de su
zona de dominio \citep{msft_activedirectory}.

\textbf{Registros m\'as relevantes para esta actividad:}
\begin{itemize}
  \item \textbf{A:} asocia un nombre de host a una direcci\'on IPv4
        (\texttt{web.sudoers.lan. IN A 192.168.0.2}).
  \item \textbf{SRV:} indica qu\'e servidor atiende qu\'e servicio en qu\'e
        puerto (\texttt{\_ldap.\_tcp} $\rightarrow$ \texttt{dc1}:389).
  \item \textbf{NS y SOA:} declaran la autoridad y la vigencia de la zona.
  \item \textbf{PTR:} asocia una direcci\'on IPv4 a un nombre (zona inversa).
  \item \textbf{CNAME:} crea un alias de otro nombre
        (\texttt{www IN CNAME web}).
\end{itemize}

\subsubsection{El DNS integrado de Samba (SAMBA\_INTERNAL)}

Samba 4 integra en un \'unico daemon los servicios de Active Directory:
LDAP, Kerberos, SMB y el \textbf{DNS} \citep{sambateam2024}. Con el backend
\textbf{SAMBA\_INTERNAL} la zona \texttt{sudoers.lan} la gestiona
autom\'aticamente el propio DC: los clientes y los administradores la
actualizan con \texttt{samba-tool dns} (o con las herramientas de AD de
Windows), y la zona se reescribe din\'amicamente sin archivos de zona
manuales como los que usa BIND 9 standalone.

\subsubsection{El servidor web y el proxy inverso}

Un servidor web atiende peticiones HTTP y entrega documentos a los
navegadores. En el proyecto se utiliza \textbf{Nginx} \citep{nginxDocs},
desplegado como contenedor Docker, con una doble funci\'on: servir el
\textbf{portal est\'atico} de la intranet y actuar como \textbf{proxy inverso}
que enruta cada nombre \texttt{*.sudoers.lan} hacia el servicio correcto. Al
proxy llegan todas las peticiones y \'el las reparte; adem\'as \textbf{termina
el TLS}: el cifrado (HTTPS) se aplica en el punto \'unico de entrada y los
servicios internos se comunican entre s\'i por HTTP.

\subsubsection{Infraestructura del laboratorio}

\begin{itemize}
  \item Servidor \'unico: \textbf{Debian 12}, hostname \texttt{dc1}, IP fija
        \texttt{192.168.0.2}.
  \item DNS del dominio: \textbf{Samba 4 AD DC} (backend SAMBA\_INTERNAL),
        puerto 53; zona autoritativa \texttt{sudoers.lan}.
  \item Servidor web: \textbf{Nginx} en contenedor Docker, puertos 80/443;
        resolver del dominio \texttt{sudoers.lan}.
  \item Red: \texttt{192.168.0.0/24}; DHCP (Kea) entrega a los clientes el
        DNS \texttt{192.168.0.2} y el dominio \texttt{sudoers.lan}.
  \item Clientes: estaciones GNU/Linux y Windows unidas al dominio que
        consultan el DNS del DC y acceden al sitio por nombre.
\end{itemize}